\section{Introduction and the main result}
In the proper hat-guessing game, the players occupy the vertices of a graph,
and the adversary assigns colors subject to the condition that adjacent
vertices receive different colors. Each player observes the colors of its
neighbors and simultaneously guesses its own color. A strategy is winning
if at least one guess is correct for every permitted assignment.
Adriaensen et al.\ introduced this variant and its parameter
$\HG(G)$ in~\cite{proper}. Their Table~1, in version~4, records
$4\leq\HG(C_5)\leq5$.
The task considered here is precisely to decide between these two values.
The four-color construction is short; the five-color exclusion is the
part that requires the finite computation.

This manuscript develops a computer-assisted argument for the following
assertion.
\begin{theorem}[Main theorem; computer-assisted]\label{thm:main}
The proper hat-guessing number of the five-cycle is
\[
\HG(C_5)=4.
\]
\end{theorem}

The nonexistence argument concerns five arbitrary local functions. It does
not assume that the players use the same rule, that left and right are
interchangeable, or that a strategy is invariant under permutations of the
vertices or colors. Symmetries are used only to select representatives of
equivalent candidates.

The principal statistic is the number of local observations on which a
fixed color is guessed. There are 125 entries in the five strategy tables.
After discarding guesses that can never be correct, at least one color is
guessed on at most 25 entries. Configurations in which this color occurs
twice force its total count to be at least 17 and impose restrictions on
the five individual counts. The resulting finite classification has
1413 representatives. Some representatives are refined using equality
patterns of five diagonal answers, producing 8613 terminal Boolean
formulas. The original local computation supplies a refutation of every
terminal formula. The subsequent complete audit replays each refutation
with a new checker and separately rechecks that the actual inputs and
branching rules preserve every candidate winning strategy.

\subsection*{Why the argument has two parts}
Fix a color and first consider only proper colorings in which it appears
twice. This restricted family makes the interaction between two strategy
tables visible and yields the lower bound 17. It does not characterize
winning strategies for the whole game: a saved strategy attains 17 on
this family but loses on colorings in the other two families. We include
that witness in Appendix~\ref{app:witness}, so the change from a local
combinatorial argument to a full Boolean model has an explicit reason.

The proof has three obligations, each with a different kind of evidence:
\begin{enumerate}
\item \emph{Coverage:} every hypothetical winning strategy can be
transformed into a listed branch. Sections~\ref{sec:local}
and~\ref{sec:classification} give the mathematical reductions.
\item \emph{Input preservation:} that strategy would satisfy the actual
branch formula after suitable auxiliary variables are assigned.
Section~\ref{sec:encoding} specifies the encoding and this extension check.
\item \emph{Exclusion:} each terminal formula has a valid refutation.
Section~\ref{sec:certificates} gives the inference rule and identifies the
recorded finite verification on which the conclusion depends.
\end{enumerate}
A refutation alone establishes only the last obligation. Conversely, a
complete list of branches does not establish that any branch is impossible.

The organization follows this logical dependence rather than the
chronological order of the experiments. Section~\ref{sec:preliminaries}
fixes the model and normalization. Section~\ref{sec:local} proves the
combinatorial count restrictions. Section~\ref{sec:classification}
establishes the complete finite classification.
Sections~\ref{sec:encoding} and~\ref{sec:certificates} explain the Boolean
encoding and certificate checks. Section~\ref{sec:conclusion} gives the
final deduction. Appendix~\ref{app:reproduction} connects the numbered
results to the organized research archive and gives a procedure for
independent reproduction.

\subsection*{Status of this manuscript}
This manuscript was prepared with computer and AI assistance. On
October 26, 2026, a new input auditor and a new RUP checker completed a
full local audit of the adopted evidence. The combinatorial arguments
were also examined, with exhaustive finite cross-checks; no error
affecting the main conclusion was found in that audit. Its precise scope
is given in Section~\ref{sec:full-audit} and Appendix~\ref{app:audit}.

Here ``independent implementation'' means separate code, not independent
human authorship or external peer review. The new auditor imports no
earlier generator or auditor, and the new checker was written before
the earlier checker sources were inspected. All implementations were
nevertheless developed within the same AI-assisted workflow. The
mathematical argument, source-to-model correspondence, and checking
software remain open to external scrutiny; no proof-assistant
verification is claimed. Bibliographic statements about~\cite{proper}
refer to the inspected version, not to an exhaustive assessment of
subsequent literature.

